\chapter*{Lời Mở Đầu}
\addcontentsline{toc}{chapter}{Lời Mở Đầu}
\markboth{Lời Mở Đầu}{Lời Mở Đầu}

\begin{truyen}{Lời Mở Đầu}{Opening Words}
\chuhoa{T}{rong rất nhiều gia đình, cha mẹ thương con nhưng nói không tới. Con cái cần được hiểu nhưng lại chọn im lặng, chống đối, hoặc trả lời cho qua. Thế là cùng ở một nhà mà vẫn có những khoảng cách dài hơn cả quãng đường đi học mỗi ngày.}
\textit{(In many families, parents love deeply but fail to speak in a way that reaches. Children need understanding but answer with silence, resistance, or avoidance. So people can live in one home and still feel far apart.)}

Quyển sách này gom lại 100 màn đối thoại rất gần với đời sống gia đình: chuyện học, chuyện điểm số, chuyện áp lực, tiền bạc, kỳ vọng, tự do, tổn thương, và những câu hỏi khó mà cả cha mẹ lẫn con cái đều không dễ trả lời.
\textit{(This book gathers 100 dialogues close to real family life: learning, grades, pressure, money, expectations, freedom, hurt, and difficult questions that neither parents nor children answer easily.)}

Quyển này không viết để kết luận bên nào đúng hơn. Quyển này viết để nhắc rằng một gia đình bền không phải gia đình không cãi nhau, mà là gia đình còn chịu ngồi lại, chịu nghe, và chịu sửa cách yêu thương của mình.
\textit{(This volume is not written to declare one side more right. It is written to remind us that a strong family is not one without conflict, but one still willing to sit down, listen, and repair the way love is expressed.)}
\end{truyen}

\begin{baihoc}
Trong một ngôi nhà, nhiều khi điều đau nhất không phải là thiếu yêu thương, mà là yêu thương không tới nơi.
\textit{(In a home, the deepest pain is often not lack of love, but love that fails to arrive.)}
\end{baihoc}
\ngancach